\section{Linearni jezik}

Programsk jezik Line Mix sestavlja sedem glavnih datotek
%
\begin{itemize}
    \item \verb|lexer.mll|. Tu definiramo katere besede in znaki so del jezika. Na sliki~\ref{cd:lexer} je primer kako izgleda del datoteke \verb|lexer.mll|. Vsak znak dobi svoje ime, ki jih uporabimo v datoteki \verb|parer.mly|, da definiramo, kako izgledajo izrazi.
\begin{figure}[h]
    \begin{verbatim}
  | '!'             { BANG }
  | '^'             { CARET }
  | '%'             { MOD }
  | '&'             { AMP }
  | '('             { LPAREN }
  | ')'             { RPAREN }
  | '*'             { TIMES }
  | '+'             { PLUS }
    \end{verbatim}
    \caption{Del kode iz \(\cd{lexer.mll}\)}\label{cd:lexer}
\end{figure}    

    \item \verb|parser.mly|. V \verb|parser.mly|  je definirano, kako izgledajo izrazi. Pomembno je, da program ve v kakšnem vrstnem redu mora evalvirati izraz. Če dobi izraz \(\cd{f a b}\) mora program to razumeti kot \(\cd{(f a) b}\) namesto \(\cd{f (a b)}\).
    Prav tako mora vedeti, da \verb|^a & b| 
    pomeni \verb|(^a) & b|, saj se ">\verb|^|"< veže močneje kot ">\verb|&|"<.
    Zato j potrebno definirati \verb|atom_expr|, kot je na sliki~\ref{cd:atom_expr}.
    %
\begin{figure}[h]
     \small
        \begin{verbatim}
    atom_expr:
      | LPAREN e = expr RPAREN { e }
      | LPAREN e1 = expr COMMA e2 = expr RPAREN { Pair (e1, e2) }
      | LPAREN e = expr COLON t = ltype RPAREN  { Annot (e, t) }
      | x = VAR     { Var x }
      | n = INT     { Int n }
      | TRUE        { Bool true }
      | FALSE       { Bool false }
      | LBRACKET es = separated_list(COMMA, expr) RBRACKET { Array es }
      | CARET e = atom_expr { Demote e }
        \end{verbatim}
        \caption{Definiranje atomksih izrazov}\label{cd:atom_expr}
\end{figure}
    % TODO: Ali naj opišem kako se nekaj drži drugega
    
    
    
    
    
    

    \item \verb|syntax.ml|. 
    \item \verb|interpret.ml|. 
    \item \verb|typechecker.ml|. 
    \item \verb|line_mix.ml|. 
\end{itemize}